\subsection{Berechnung der Messunsicherheiten}
\label{sec:Messunsicherheiten}

Zur Berechnung des Mittelwertes einer Größe $x$ mit $N$ Messwerten wird die Formel
%
\begin{equation}
    \bar{x} = \frac{1}{N} \sum_{i=1}^N x_i
\end{equation}
%
benutzt. Die zugehörige Standardabweichung des Mittelwertes $\increment \bar{x}$ berechnet sich über
%
\begin{equation}
    \increment \bar{x} = \sqrt{\frac{1}{N \l(N-1\r)}\sum_{i=1}^N \left(x_i - \bar{x}\right)^2}
\end{equation}
%
Die Unsicherheit einer zusammengesetzen Größe $g\left(x_1, …, x_n\right)$ wird mithilfe der Gaußschen
Fehlerfortpflanzung bestimmt
%
\begin{equation}
    \increment g\left(x_1, …, x_n\right) = \sqrt{\sum_{i=1}^N \left(\increment x_i ~ \partial_i \, f \right)^2}
\end{equation}
%
Größen werden in der Form $\bar{x} \pm \increment \bar{x}$ angegeben. Alle Berechnungen, Graphen, sowie die 
Bestimmung der Messunsicherheiten werden mit Python 3.11.5 und entsprechenden Bibliotheken
\footnote{Numpy \cite{numpy}, Matplotlib \cite{matplotlib} und Uncertainties \cite{uncertainties}} durchgeführt.